\documentclass[10pt,a4paper]{amsart}
\usepackage[margin=19mm]{geometry}
\usepackage{amsmath,amssymb,mathtools}
\usepackage[scr=rsfs,cal=euler]{mathalfa}
\usepackage{xfrac}
\usepackage[unicode,hidelinks,bookmarks=false]{hyperref}
\newcommand{\ie}{i.e.\ }
\newcommand{\cf}{cf.\ }
\newcommand{\resp}{respectively\ }
\newcommand{\Subsection}{\subsection}
\providecommand{\nobreakdash}{\nobreak\hbox{-}}
\setlength{\emergencystretch}{2em}
\multlinegap0pt
\usepackage{bm}
\usepackage{amssymb}
\usepackage{algorithm}\usepackage{algpseudocode}
\usepackage[shortlabels]{enumitem}
\newtheorem{proposition}{Proposition}
\newtheorem{lemma}{Lemma}
\newcommand{\Op}{\operatorname{Op}}
\newcommand{\dd}{\,\mathrm{d}}
\newcommand{\ii}{\mathrm{i}}
\title{Time-modulated refocusing in inhomogeneous medium}
\author{Hongyu Liu$^\P$, Wei Wu$^\dag$}
\date{Source: arXiv:2608.29542v1 (CC BY 4.0)}
\begin{document}
\maketitle

\noindent\emph{Source and licence.} Time-modulated refocusing in inhomogeneous medium. Version arXiv:2608.29542v1 (CC BY 4.0), distributed by arXiv under the Creative Commons Attribution 4.0 International licence, http://creativecommons.org/licenses/by/4.0/, as recorded in the arXiv licence metadata for this exact version and shown on the abstract page https://arxiv.org/abs/2608.29542v1. Full licence notice: \texttt{licenses/arxiv-2608.29542-CC-BY-4.0.txt}.\par\smallskip

\noindent\textit{Editorial scope.} Counted unit: the article's lemma eq:measured-energy (source0.tex lines 1509--1523). The article's complete original proof of it is delivered with it (source0.tex lines 1525--1646, one \texttt{proof} environment of 122 lines). No mathematics is written, repaired or re-derived: the statement and the proof are the article's own lines, in the article's own notation, and no retained line is substituted or re-flowed.  Retained uncounted as the source-local context the proof needs: lines 484--488; lines 553--555; lines 589--598; lines 616--618; lines 790--795; lines 821--823; lines 876--878; lines 1325--1330; lines 1445--1456; lines 1494--1505.  Nothing was omitted from any retained span, so the package carries no omission record.  No retained line contains a citation, so the wrapper carries no bibliography and no bibliography entry.  No retained line contains a figure environment, an image-inclusion command or drawing code, so no image is delivered and the package declares no asset dependency.  Editorial formatting only, in addition: an AMS \texttt{amsart} preamble that restates the article's own declarations and macros used by the retained lines, namely the environments \texttt{proposition}, \texttt{lemma} on separate counters, together with the standard packages those lines need.  The retained environments print in this document as Proposition 1, Lemma 1 in order of appearance, whereas the article prints the same items with the numbering of its own full document.  The complete native source of the article as distributed in the arXiv e-print for this exact version is bundled unchanged as \texttt{sources/208863/source0.tex}.
\begin{equation}\label{eq:explicit-r-positive-branch}
    \bm r(\bm x,\bm\xi)
    =\left(\frac{\widehat\xi_2}{\sqrt{2\mu}}, -\frac{\widehat\xi_1}{\sqrt{2\mu}}, \frac{1}{\sqrt{2\epsilon_0\epsilon_\infty(\bm x)}}, 0\right)^\top, \qquad
    \widehat{\bm\xi}:=\frac{\bm\xi}{|\bm\xi|},
\end{equation}

\begin{equation}\label{eq:semiclassical-system}
    \left(\partial_t+\frac{\ii}{h}\Op_h(A_0^{-1}\widetilde\Lambda)+\Op_h(A_0^{-1}B_0)\right)\bm W=0.
\end{equation}

\begin{equation}\label{eq:symbol-product}
    a_h\# b_h = a_hb_h+\frac{h}{\ii}\sum_{j=1}^2(\partial_{\xi_j}a_h)(\partial_{x_j}b_h)+h^2r_{a,b,h},
\end{equation}
where $r_{a,b,h}$ has order two lower than the product $a_hb_h$.

\subsubsection{Construction of injection symbol.}\label{sec:injection-symbol}
Suppose that a scalar amplitude $u$ is inserted into the vector system through a symbol $\bm r_h$. We want to find a symbol $\bm r_h$, such that a scalar amplitude $u$, applied by $\Op_h(\bm{r}_h)$, is the vector solution of \eqref{eq:semiclassical-system}, and 
\begin{equation}\label{eq:intertwine}
\left(\partial_t+\frac{\ii}{h}\Op_h(\mathsf M_{\rm md})+\Op_h(\mathsf C_{\rm md})\right)\Op_h(\bm r_h)u=\Op_h(\bm r_h)\left(\partial_t+\frac{\ii}{h}\Op_h(\omega)+\Op_h(\beta_h)\right)u \quad\text{modulo }O(h^\infty).
\end{equation}

\begin{equation}\label{eq:first-right-correction}
    (\mathsf M_{\rm md}-\omega)\bm r_1 = \ii\left(\bm{\mathcal D}+\mathsf C_{\rm md}\bm r-\bm r\beta\right).
\end{equation}

\begin{align}
    U_0(t)J_h&=J_hS_h(t)+\mathcal R_h^{\infty,+}(t),
    \label{eq:intertwining-operator-remainder}\\
    J_h^{\#}U_0(t)&=S_h(t)J_h^{\#}+\widetilde{\mathcal R}_h^{\infty,+}(t),
    \label{eq:left-intertwining-operator-remainder}
\end{align}

\begin{equation}\label{eq:egorov-with-class}
    F_h(-t)\Op_h(a_h)F_h(t)=\Op_h(a_h\circ\bm\Phi^t)+hR_{a,h}(t),
\end{equation}

\begin{equation}\label{eq:damped-factorization}
    S_h(t)=\Op_h(e^{-\mathcal G_t})F_h(t)+h\mathcal R_h^{\rm damp}(t),
\end{equation}

\begin{equation}\label{eq:source-state-before-time-derivative}
 Q_{\star,h}J_h^{\#}\bm W_h(0)
 =\frac1{(2\pi h)^2}\int e^{\ii(\bm x-\bm z_\star)\cdot\bm\xi/h}
 \left[s_h\bigl(\omega(\bm z_\star,\bm\xi)\bigr)\chi_\star(\bm\xi)
 +h a_{0,h}(\bm x,\bm\xi)\right]\dd\bm\xi,
\end{equation}

\begin{proposition}\label{prop:inverse-oa}
On a neighborhood of the selected measured component there exists a uniformly bounded order-zero semiclassical pseudodifferential operator $\mathcal L_h^+$ such that
\begin{equation}\label{eq:observation-inverse}
    \mathcal L_h^+\mathcal O_{\mathcal A}^+J_h
    = I+hR_h^{\rm inv},
\end{equation}
microlocally on the selected component. The operator $R_h^{\rm inv}$ is uniformly bounded on $\mathcal H_{\rm s}$ with
\begin{equation*}
 \|hR_h^{\rm inv}\|_{\mathcal H_{\rm s}\to\mathcal H_{\rm s}}\le Ch.
\end{equation*}
The full symbol of $\mathcal L_h^+$ may be chosen compactly supported with uniform order-zero derivative bounds, and its principal symbol is $o_+^{-1}$ on the selected component.
\end{proposition}

\begin{equation}\label{eq:measured-coefficient}
 \begin{aligned}
 \mathscr P_h(\varpi,\theta)
 &:=\left\langle
 Q_{\star,h}F_h(-t_{0,h})\mathcal L_h^+d_h(t_{0,h}),
 \psi_{\bm\Psi_{T_h}(\varpi,\theta),h}\right\rangle,\\
 \mathcal H_h(\varpi,\theta)
 &:=\exp\left[2\int_0^{t_{0,h}}
 \kappa\bigl(\bm\Phi^s(\bm\Psi_{T_h}(\varpi,\theta))\bigr)\dd s\right]
 |\mathscr P_h(\varpi,\theta)|^2.
 \end{aligned}
\end{equation}

\begin{lemma}
For every $N\in\mathbb N_0$, there is a function $r_{E,h}$ such that
\begin{equation}\label{eq:measured-energy}
 \mathcal H_h(\varpi,\theta)
 =\frac{1}{\pi h}|s_h(\varpi)|^2
 \exp\left[-2\int_0^{t_{0,h}}
 \gamma\bigl(\bm\Phi^s(\bm\Psi_{T_h}(\varpi,\theta))\bigr)\dd s\right]
 \bigl(1+h r_{E,h}(\varpi,\theta)\bigr),
\end{equation}
and
\begin{equation*}
 \|r_{E,h}\|_{C^N(I_{\star,T_h}\times\Theta_{T_h})}\le C_N
\end{equation*}
uniformly for small $h$.
\end{lemma}

\begin{proof}
We start with studying $Q_{\star,h}F_h(-t_{0,h})\mathcal L_h^+d_h(t_{0,h})$ in the definition of $\mathscr P_h$. By the definition of $d_h(t_{0,h})$, we have
\begin{align*}
    \mathcal L_h^+d_h(t_{0,h}) = \mathcal L_h^+\mathcal O_{\mathcal A}^+U_0(t_{0,h})\bm{W}_h(0).
\end{align*}
Now define $\widetilde{\chi}_\star$ satisfying $\widetilde{\chi}_\star=1$ on $\operatorname{supp}(\chi_\star)$, and set $\widetilde Q_{\star,h}:=\Op_h(\widetilde\chi_\star)$. We want to prove 
\begin{equation}\label{eq:theorem-proof-1}
    \mathcal L_h^+\mathcal O_{\mathcal A}^+U_0(t_{0,h})(I-J_h\widetilde{Q}_{\star,h}J_h^\#)\bm{W}_h(0)
\end{equation}
is negligible. Notice that from \eqref{eq:intertwining-operator-remainder} and \eqref{eq:left-intertwining-operator-remainder},
$$
    U_0(t)(I-J_hJ_h^\#) = (I-J_hJ_h^\#)U_0(t)+R_h^\infty(t).
$$
Rewrite \eqref{eq:theorem-proof-1} as
$$
    \mathcal L_h^+\mathcal O_{\mathcal A}^+U_0(t_{0,h})\left((I-J_hJ_h^\#) + J_h(I-\widetilde{Q}_{\star,h})J_h^\#\right).
$$
On one hand, recall that in the proof of Proposition \ref{prop:inverse-oa}, $\mathcal L_h^+=\Op_h(\chi/o_+)$ with $\chi\equiv 1$ on the support of $\chi_\star$, so 
$$
    \mathcal L_h^+\mathcal O_{\mathcal A}^+U_0(t_{0,h})J_h(I-\widetilde{Q}_{\star,h})J_h^\#\bm{W}_h(0)
$$
is negligible. On the other hand, from proof of Proposition \ref{prop:inverse-oa}, the support of $\chi$ does not meet any other observed sheets, hence $\mathcal L_h^+\mathcal O_{\mathcal A}^+$ applying to $I-J_hJ_h^\#$ is also negligible, which means that
$$
    \mathcal L_h^+\mathcal O_{\mathcal A}^+U_0(t_{0,h})(I-J_hJ_h^\#)\bm{W}_h(0)
$$
is negligible. This finishes the proof of the claim that \eqref{eq:theorem-proof-1} is negligible. We then have
\begin{align*}
    \mathcal L_h^+d_h(t_{0,h}) &= \mathcal L_h^+\mathcal O_{\mathcal A}^+U_0(t_{0,h})J_h\widetilde{Q}_{\star,h}J_h^\#\bm{W}_h(0)+R_h^\infty\bm{W}_h(0)\\ 
    &= \mathcal L_h^+\mathcal O_{\mathcal A}^+J_hS_h(t_{0,h})\widetilde{Q}_{\star,h}J_h^\#\bm{W}_h(0)+R_h^\infty\bm{W}_h(0)\\
    &= (I+hR_h^{\rm inv})S_h(t_{0,h})\widetilde{Q}_{\star,h}J_h^\#\bm{W}_h(0)+R_h^\infty\bm{W}_h(0)
\end{align*}
Applying $Q_{\star,h}F_h(-t_{0,h})$ from the left yields
\begin{align*}
    &Q_{\star,h}F_h(-t_{0,h})\mathcal L_h^+d_h(t_{0,h}) \\ 
    &=Q_{\star,h}F_h(-t_{0,h})S_h(t_{0,h})\widetilde{Q}_{\star,h}J_h^\#\bm{W}_h(0)+hQ_{\star,h}F_h(-t_{0,h})R_h^{\rm inv}S_h(t_{0,h})\widetilde{Q}_{\star,h}J_h^\#\bm{W}_h(0)+R_h^\infty\bm{W}_h(0).
\end{align*}
$R_h^\infty$ is negligible. Define $R_h:=Q_{\star,h}F_h(-t_{0,h})R_h^{\rm inv}S_h(t_{0,h})\widetilde{Q}_{\star,h}$. $R_h$ is uniformly of order 0 from the property of $F_h$, $R_h^{\rm inv}$ and $S_h$. Truncation operators $Q_{\star,h}$ and $\widetilde{Q}_{\star,h}$ don't change the order. Then 
\begin{equation}\label{eq:measured-data-reduction}
    Q_{\star,h}F_h(-t_{0,h})\mathcal L_h^+d_h(t_{0,h}) =Q_{\star,h}F_h(-t_{0,h})S_h(t_{0,h})\widetilde{Q}_{\star,h}J_h^\#\bm{W}_h(0)+hR_hJ_h^\#\bm{W}_h(0)+R_h^\infty\bm{W}_h(0).
\end{equation}

Now we calculate the leading order part in \eqref{eq:measured-data-reduction}. We first write $\kappa$ in terms of the known principal symbol $\omega$. From \eqref{eq:explicit-r-positive-branch},
\begin{equation*}
 r_3(\bm x,\bm\xi)=\frac{1}{\sqrt{2\epsilon_0\epsilon_\infty(\bm x)}}.
\end{equation*}
Since the entries of $\bm r$ are real, $A_0$ is independent of $\bm\xi$, and $\bm r^*A_0\bm r=1$, differentiation in $\xi_j$ gives
\begin{equation*}
 \bm\ell\partial_{\xi_j}\bm r=0.
\end{equation*}
The first two components of $\bm r$ are independent of $\bm x$. Recall $M(\bm x,\bm\xi) = A_0^{-1}(\bm x)\widetilde\Lambda(\bm\xi)$. We have
\begin{equation*}
 \sum_{j=1}^2(\partial_{\xi_j}M)(\partial_{x_j}\bm r)=\frac1\mu(\partial_{x_2}r_3,-\partial_{x_1}r_3,0,0)^\top.
\end{equation*}
Using
\begin{equation*}
 \bm\ell=\left(\sqrt{\frac\mu2}\widehat\xi_2,-\sqrt{\frac\mu2}\widehat\xi_1,\sqrt{\frac{\epsilon_0\epsilon_\infty}2},0\right)
\end{equation*}
the definition of $\kappa$ in section \ref{sec:injection-symbol} and the definition of $\bm{\mathcal D}$ preceding \eqref{eq:first-right-correction}, we obtain
\begin{equation*}
 \kappa(\bm x,\bm\xi)=\bm\ell\bm{\mathcal D}(\bm x,\bm\xi)=\frac12\widehat{\bm\xi}\cdot\nabla_{\bm x}\bigl(\mu\epsilon_0\epsilon_\infty(\bm x)\bigr)^{-1/2}=\frac12\sum_{j=1}^2
 \partial_{x_j}\partial_{\xi_j}\omega(\bm x,\bm\xi).
\end{equation*}

We next calculate the local coefficient at $t=0$. Write $\bm\Psi_{T_h}(\varpi,\theta)=(\bm z_\star,\bm\xi_0)$. On the retained chart, $\chi_\star=1$ near $\bm\xi_0$. From \eqref{eq:source-state-before-time-derivative}, direct integration in $\bm x$ gives
\begin{equation*}
 \int_{\mathbb R^2}e^{\ii(\bm x-\bm z_\star)\cdot(\bm\xi-\bm\xi_0)/h}e^{-|\bm x-\bm z_\star|^2/(2h)}\dd\bm x=2\pi h e^{-|\bm\xi-\bm\xi_0|^2/(2h)}.
\end{equation*}
Consequently,
\begin{equation}\label{eq:source-local-coefficient}
 \begin{aligned}
 &\left\langle \widetilde{Q}_{\star,h}J_h^{\#}\bm W_h(0),
 \psi_{\bm\Psi_{T_h}(\varpi,\theta),h}\right\rangle\\
 &=(2\pi h)^{-1}(\pi h)^{-1/2}
 \int e^{-|\bm\xi-\bm\xi_0|^2/(2h)}
 \left[s_h\bigl(\omega(\bm z_\star,\bm\xi)\bigr)
 +h\widetilde a_{0,h}(\bm\xi)\right]\dd\bm\xi\\
 &=(\pi h)^{-1/2}s_h(\varpi)
 \bigl(1+h r_{0,h}(\varpi,\theta)\bigr).
 \end{aligned}
\end{equation}
The last equality follows by setting $\bm\xi=\bm\xi_0+h^{1/2}\bm\zeta$ and expanding the amplitude at $\bm\xi_0$. The terms of order $h^{1/2}$ are odd in $\bm\zeta$, so
$$
    \int e^{-|\bm\zeta|^2/2}\zeta_j\dd\bm\zeta=0,    
$$  
and the remaining error is $O(h)$. After differentiating in $(\varpi,\theta)$, each differentiated integrand is bounded by a polynomial in $\bm\zeta e^{-|\bm\zeta|^2/2}$, and hence is integrable uniformly in $h$. Since $|s_h|$ is bounded below on $I_{\star,T_h}$, the error can be written in the relative form used in \eqref{eq:source-local-coefficient}.

Since 
\begin{equation*}
 \mathcal G_t(\bm\Phi^t(\bm\rho))=\int_0^t\beta(\bm\Phi^{t-s}(\bm\rho))\dd s=\int_0^t\beta(\bm\Phi^s(\bm\rho))\dd s,
\end{equation*}
from \eqref{eq:damped-factorization} and \eqref{eq:egorov-with-class},
\begin{equation}\label{eq:interaction-propagator}
 F_h(-t)S_h(t)=\Op_h\left(\exp\left[-\int_0^t\beta\bigl(\bm\Phi^s(\cdot)\bigr)\dd s\right]\right)+hR_h(t)
\end{equation}
for $0\le t\le t_+$ on the selected source-side set, where $R_h(t)$ is uniformly of order zero. 

Apply \eqref{eq:interaction-propagator} to the oscillatory-integral representation of $\widetilde{Q}_{\star,h}J_h^{\#}\bm W_h(0)$ in \eqref{eq:source-state-before-time-derivative}. The left symbol product first gives the principal amplitude
\begin{equation*}
 \exp\left[-\int_0^t\beta\bigl(\bm\Phi^s(\bm x,\bm\xi)\bigr)\dd s\right]s_h\bigl(\omega(\bm z_\star,\bm\xi)\bigr).
\end{equation*}
Taylor expansion in $\bm x$ about $\bm z_\star$, followed by
\begin{equation*}
 (x_j-z_{\star,j})e^{\ii(\bm x-\bm z_\star)\cdot\bm\xi/h}
 =\frac h\ii\partial_{\xi_j}e^{\ii(\bm x-\bm z_\star)\cdot\bm\xi/h},
\end{equation*}
shows that replacing $\bm x$ by $\bm z_\star$ changes the amplitude by $h$ times a controlled amplitude. Thus, uniformly for $0\le t\le t_+$,
\begin{equation}\label{eq:processed-field}
 \begin{aligned}
 &Q_{\star,h}F_h(-t)S_h(t)\widetilde{Q}_{\star,h}J_h^{\#}\bm W_h(0)\\
 &\quad=\frac1{(2\pi h)^2}\int e^{\ii(\bm x-\bm z_\star)\cdot\bm\xi/h}\bigg[\exp\left(-\int_0^t\beta\bigl(\bm\Phi^s(\bm z_\star,\bm\xi)\bigr)\dd s\right)
 s_h\bigl(\omega(\bm z_\star,\bm\xi)\bigr)+h a_{t,h}(\bm x,\bm\xi)\bigg]\dd\bm\xi,
 \end{aligned}
\end{equation}
where $a_{t,h}$ is a controlled amplitude with uniform bounds after differentiation in $t$.

Pairing \eqref{eq:processed-field} at $t=t_{0,h}$ with the test function in \eqref{eq:measured-coefficient}, and then using \eqref{eq:measured-data-reduction}, gives
\begin{equation}\label{eq:processed-coefficient}
 \mathscr P_h(\varpi,\theta)=(\pi h)^{-1/2}s_h(\varpi)\exp\left[-\int_0^{t_{0,h}}\beta\bigl(\bm\Phi^s(\bm\Psi_{T_h}(\varpi,\theta))\bigr)\dd s\right]
 \bigl(1+h r_{P,h}(\varpi,\theta)\bigr),
\end{equation}
where $r_{P,h}$ is uniformly bounded in every $C^N$ norm. Finally, $\beta=\gamma+\kappa$. Multiplying the square of \eqref{eq:processed-coefficient} by the exponential factor in \eqref{eq:measured-coefficient} cancels the term containing $\kappa$ and gives \eqref{eq:measured-energy}.
\end{proof}

\end{document}
